% =====================================================================
%  5 — PLAN DE LA SOUTENANCE : les trois axes
% =====================================================================
% Deux parties : (1) les enjeux sociaux et de productivité qui motivent la
% réduction de la dette technique (dette technique, intérêt de l'IA, gains
% de productivité) ; (2) le cœur du propos, écrire du code de production
% dans une base déjà fortement endettée. Mise en page : gros numéros serif
% en accent, titres serif encre, descriptif en gris, filet de séparation.
\newcommand{\planitem}[3]{%
  \noindent
  \begin{minipage}[c]{1.85cm}%
    {\seriffont\bfseries\fontsize{38}{38}\selectfont\color{accent}#1}%
  \end{minipage}%
  \begin{minipage}[c]{\dimexpr\textwidth-1.95cm\relax}%
    {\seriffont\bfseries\fontsize{16}{20}\selectfont\color{ink}#2\par}%
    \vspace{2pt}%
    {\fontsize{10.5}{14}\selectfont\color{ink2}#3\par}%
  \end{minipage}\par}

\begin{frame}{Trois points qui me tiennent à cœur}
  \vspace{0.30cm}
  \planitem{01}{Les enjeux sociaux de la dette technique}%
    {Motivation des équipes, dette de connaissance et coût humain du legacy}

  \vspace{0.28cm}{\color{rule}\rule{\textwidth}{0.6pt}}\par\vspace{0.28cm}

  \planitem{02}{L'IA, une solution pleine de promesses}%
    {Un levier puissant mais conditionné, et ses propres risques de dette}

  \vspace{0.28cm}{\color{rule}\rule{\textwidth}{0.6pt}}\par\vspace{0.28cm}

  \planitem{03}{Écrire du code dans un existant fortement endetté}%
    {L'adressage incrémental et l'isolation du code neuf sur une base legacy}
\end{frame}
